\chapter{Modification apportée au software (DSP)}

\begin{enumerate}
    \item Machine Etat
    \item Conversion courant et pos modif par rapport a
    \item Structure et propeté du code
    \item correction de la pos
\end{enumerate}

\section{Introduction}



La mise en place de Code Composer Studio est disponible dans le manuel \autoref{sec:MiseEnPlace}.

\section{Modifications mineur}
\subsection{Amélioration du code via l'IA}
Afin d'optimiser l'appropriation du code source hérité, une première révision syntaxique a été réalisée à l'aide d'outils d'intelligence artificielle. Cette démarche avait pour objectif d'améliorer la structure du programme, d'enrichir les commentaires et d'en accroître la lisibilité globale. Le comportement de la maquette a été rigoureusement contrôlé à la suite de cette restructuration purement formelle. Les résultats obtenus étant strictement identiques à ceux de la version d'origine, il est confirmé que cette intervention n'a introduit aucun dysfonctionnement supplémentaire.

\subsection{Création d'un fichier secondaire .c}
Ce fichier constitue une révision du travail initial mené par Thomas Freyche. Afin d'améliorer la clarté et la modularité du code source, les définitions ont été scindées sur deux fichiers distincts, bien que la logique de fonctionnement demeure strictement identique.

\subsection{Modification des coefficient du coupen bande}
Toutefois, lors de l'analyse du code source hérité, il a été constaté que ces définitions n'étaient pas utilisées. Une autre variante de coefficients était active :

\begin{minted}{c}
//bandstop Filter 56Hz 2nd order bande coupe de 47Hz  90Hz (0) (test)
#define A1_B1_0         -1.980143772501098
#define A2_0            0.980339665209542
#define B0_B2_0         0.9901698326047711
\end{minted}

Cette disparité provient vraisemblablement d'essais expérimentaux antérieurs non documentés. Par souci de cohérence avec le travail de Laucella \cite{Laucella}, les paramètres ciblant la fréquence de 75~Hz ont été restaurés.

\subsection{Modification de l'emplacement du filtre coupe-bande}

Afin d'éviter un retard d'une période de 40 us, le filtre fut placé juste après le calcul de la force. Précédement celui-ci se trouvais avant ce calcul et fut employé pour le calcul de la consigne de courant.

\section{Machine d'état}

\section{Regulation de courant}

\section{Regulation de position}

\subsection{Methode inverse}




